%! Author = chtompki
%! Date = 1/14/26
\documentclass[12pt]{amsart}
% Default margins are too wide all the way around. I reset them here
\setlength{\hoffset}{-.75in}
\setlength{\textwidth}{6.5in}
\setlength{\voffset}{-.5in}
\setlength{\textheight}{9.0in}
\usepackage{amsmath}
\usepackage{amssymb}
\usepackage{amsthm}
\usepackage{enumitem}
\usepackage{setspace}
\usepackage{graphicx}
\usepackage[colorlinks=true, urlcolor=black, linkcolor=black]{hyperref}

\newenvironment{solution}
{\begin{proof}[Solution]}
{\end{proof}}

\title{Numerical Analysis\\Guided Lab 2}
\author{Rob Tompkins}
\date{\today}

\begin{document}
    \maketitle
    \date{\today}
    \begin{onehalfspace}
        Note, all the files associated with Guided Lab 2 are checked into the : \\
        \underline{https://github.com/chtompki/MathSpring2026/tree/main/NumericalAnalysis/GuidedLab2}

        % Problem 1
        \textbf{Exercise 1.} We've added the code files to github at the following links:
        \begin{itemize}
            \item \underline{https://github.com/chtompki/MathSpring2026/tree/main/NumericalAnalysis/GuidedLab2}
            \item \underline{https://github.com/chtompki/MathSpring2026/tree/main/NumericalAnalysis/GuidedLab2}
        \end{itemize}
        We went about completing this problem by writing the following code:
        \begin{tiny}
        \begin{verbatim}
function Q = gram_schmidt( X )
    % Q = gram_schmidt( X )
    % X is an m by n matrix
    % Q is an m by r matrix with 1 <= r <= n that is linearly independent
    matrix_size = size(X);
    m = matrix_size(1,1);
    n = matrix_size(1,2);
    if X == zeros(m,n)
        error('There does not exist any type of basis for the zero vector space.');
    elseif n == 1
        Q = X(1:m,1)/norm(X(1:m,1));
    else
        if rank(X) ~= n
            X = basis_col(X);
        end
        matrix_size = size(X);
        m = matrix_size(1,1);
        n = matrix_size(1,2);
        Q = X(1:m,1)/norm(X(1:m,1));
        for i = 2:n
            u = X(1:m,i);
            v = zeros(m,1);
            for j = 1:(i - 1)
                v = v - dot(u,Q(1:m,j))*Q(1:m,j);
            end
            v_ = u + v;
            Q(1:m,i) = v_/norm(v_);
        end
    end
end
        \end{verbatim}
        \end{tiny}

        % Problem 2
        \textbf{2.}

        % Problem 3
        \textbf{3.}

        % Problem 4
        \textbf{4.}


    \end{onehalfspace}
\end{document}
